\documentclass[10pt,a4paper]{article}

\usepackage[T1]{fontenc}
\usepackage{newtxtext,newtxmath}
\usepackage{amsmath}
\usepackage{booktabs,tabularx,array}
\usepackage[margin=1.75cm,top=1.6cm,bottom=1.8cm]{geometry}
\usepackage{microtype}
\usepackage[font=small,labelfont=bf]{caption}
\usepackage{titlesec}
\usepackage{enumitem}
\usepackage{multicol}
\usepackage[colorlinks=true,linkcolor=blue!50!black,citecolor=green!35!black,urlcolor=blue!50!black]{hyperref}
\usepackage{xcolor}

\titleformat{\section}{\normalsize\bfseries}{\thesection.}{0.5em}{}
\titlespacing*{\section}{0pt}{1.0ex plus 0.3ex}{0.5ex}
\titleformat{\paragraph}[runin]{\normalsize\bfseries}{}{0pt}{}[.]
\titlespacing*{\paragraph}{0pt}{0.6ex plus 0.2ex}{0.6em}
\setlength{\parskip}{0pt}
\setlength{\parindent}{1.2em}
\setlist{leftmargin=1.4em,itemsep=1pt,topsep=2pt,parsep=0pt}

\newcommand{\Z}{\mathbb{Z}}
\newcommand{\Oc}{\mathcal{O}}
\newcommand{\ind}{\operatorname{ind}}
\newcommand{\Xa}{X_{7862}}
\newcommand{\Xb}{X_{7447}}

\begin{document}

\begin{center}
{\large\bfseries Search for heterotic line-bundle Standard Models}\\[3pt]
{\normalsize Summary of results, the algorithm, and directions beyond the heterotic string}\\[4pt]
{}
\end{center}

\vspace{-2pt}
{\small\noindent\textbf{Abstract.}
Much of what is known about string compactifications with the particle content
of the supersymmetric Standard Model comes from computer scans. What a scan
establishes depends on the status of each filter it applies. A filter that is
not a genuine necessary condition can throw away real models, and one that is
weaker than it could be lets spurious ones through. We revisit one of the most
thoroughly scanned cases, the $E_8\times E_8$ heterotic string on favourable
complete-intersection Calabi--Yau threefolds with a gauge bundle built from five
line bundles, and rebuild the search so that a candidate is discarded only when
a proof requires it. The main step is to replace the finite slope-stability test
of earlier scans by its complete form, which asks whether a linear space meets
a fixed cone. That question can be settled exactly on a partial model, so whole
branches of the search are cut at once. The engine reproduces the published
catalogs of the two benchmark manifolds model for model, and across all 68
manifolds it never misses a published model inside its search range. It also
finds a stable GUT model that the catalog lacks, and eleven new three-family
models on a quotient of the tetraquadric, each lying outside limits that the
earlier standard-model scan set for itself. We then ask whether any of these
models has a realistic superpotential, computing the leading couplings directly
instead of inferring them from charges. For this quotient the answer is no: no
model, published or new, combines a leading-order top-quark Yukawa coupling,
Yukawa matrices that can reach full rank, and the absence of dangerous baryon-
and lepton-number-violating operators. We end by describing how the same
approach could be used for Type IIB orientifolds, flux vacua, F-theory and
intersecting branes.\par}

\section{The problem}

A bundle $V=\bigoplus_{a=1}^5\Oc_X(k_a)$ on a favourable CICY $X$ with
$h=h^{1,1}(X)$ is an integer $h\times5$ matrix $K=(k_1,\dots,k_5)$. For $V$ to
give a supersymmetric $SU(5)$ GUT with three chiral families that descends to
an MSSM-like theory on $X/\Gamma$, with $\Gamma$ acting freely, one needs
$c_1(V)=0$, $\ind V=-3|\Gamma|$, the Green--Schwarz condition
$c_2(TX)-c_2(V)\ge0$, and poly-stability. For line bundles poly-stability means
that the five slopes $\mu_{k_a}(t)=d_{ijl}\,k_a^i t^j t^l$ vanish at a common
point of the Kähler cone. The descent then requires an equivariant structure, a
Wilson line and the $\Gamma$-representation on bundle cohomology. Earlier work
turned this into scans over boxes $|k_a^i|\le K$
\cite{AGLP11,AGLP12,ACGLP13}, producing a catalog of 63,325 GUT models on the
68 favourable CICYs and a database of standard models with their
superpotential selection rules.

Those scans leave three questions open: what a negative result proves, when
several filters are finite truncations of conditions on infinitely many charge
combinations; whether the downstairs spectrum can be computed from the group
action rather than looked up; and what the superpotential does beyond charge
conservation, which is necessary for a coupling but not sufficient.

\section{The algorithm and what is new in it}

The search is a depth-first traversal of sorted column prefixes
$k_1\preceq\dots\preceq k_m$, $m\le4$. The fifth column is fixed by
$\sum_ak_a=0$, so every class of bundles is visited exactly once. Twelve rules
prune the tree; each is a proven necessary condition, except one that is a
declared restriction of the domain, and none is a heuristic.

\paragraph{A complete slope obstruction} Earlier scans reject $K$ when some
combination $\sum_ac_ak_a$ with $c\in\{-1,0,1\}^5$ has a one-signed matrix
$M(k)_{jl}=d_{ijl}k^i$, since a one-signed quadratic form cannot vanish inside
the Kähler cone. We show that the same test over all real $c$ is a geometric
statement: the span $S$ of the charge vectors meets a fixed rational polyhedral
cone $P=\{k:M(k)\ge0\}$ away from $M=0$. Three properties make this useful in a
search. It is \emph{monotone}, so a failure on a prefix prunes every
completion. It is \emph{decided exactly} by Gordan's alternative applied to the
extreme rays of $P$ projected onto $S^\perp$, whose dimension shrinks as the
prefix grows; at the deepest and most numerous nodes it reduces to a sign check
or an angular-gap check. And it is \emph{strictly stronger} than the finite
test: on the tetraquadric $\Xa$ it removes one orbit that the finite test keeps
at $K=3$ and $K=4$, and removing it is what brings the search into exact
agreement with the published catalog.

\paragraph{Searches without a charge cutoff} A common slope zero forces a
primitive integer vector $p$ with all $p_i>0$ and $\langle k_a,p\rangle=0$ for
every $a$, so each model lies in a slice $p^\perp$. When
$(\sum_ip_i)^2>(h-1)|p|^2$, the anomaly condition bounds every charge in the
slice, and the slice can be searched over all integers. On $\Xa$, 425 of 469
such slices are proven empty for arbitrarily large charges. Separately, the
greatest common divisor $d_X$ of all line-bundle indices on $X$ must divide
$3|\Gamma|$; this excludes some pairs $(X,\Gamma)$ for any charges, and both
pairs it excludes among the 68 CICYs are indeed empty in the catalog.

\paragraph{An exact decision of stability} A bundle can pass every sign test
while its five slopes have no simultaneous zero. When the charge matrix has
rank $h-1$ we decide the question exactly: a lexicographic Gröbner basis,
isolation of the real roots of each irreducible factor of the eliminant in
rational intervals, interval arithmetic for positivity, and a check that every
slope vanishes modulo the minimal polynomial. It returns either a
real-algebraic Kähler class where $V$ is poly-stable or a proof that none
exists.

\paragraph{Selection rules as network flows} Whether the $U(1)^4$ charges allow
a superpotential monomial after any number of singlet insertions is a monoid
membership question. We show that it is a nonnegative-flow problem on a small
directed graph. Total unimodularity makes the real relaxation exact, so the
question is decided to all orders in the singlets, and a min-cost flow gives
the lowest order at which a coupling can appear.

\paragraph{Cohomology and couplings from first principles} A nonzero
$r\times r$ minor modulo a single prime certifies that a Koszul map has rank at
least $r$ for generic complex structure, with no good-reduction hypothesis. The
$\Gamma$-representation on cohomology is computed from the group action through
the Koszul resolution, Bott's theorem and equivariant map ranks. Leading-order
Yukawa couplings are computed as cup products
$H^1\times H^1\times H^1\to H^3(X,\Oc_X)$ \cite{BBCL16} after splitting each
class by Koszul level, which detects couplings that the charges allow but that
vanish for every complex structure.

\smallskip
The pruning keeps the search small. On $\Xa$, from $K=1$ to $K=7$ the nominal
domain grows by a factor $1.5\times10^{11}$ while the number of complete
bundles the search builds grows by a factor of about 700.

\section{Results}

\paragraph{Exact agreement with the published catalogs} The survivor sets
coincide with the catalog model for model: all 10 and all 58 models on $\Xa$
($\Gamma=\Z_2$ and $\Z_2\times\Z_2$) and all 93 on $\Xb$, whose counts we
predicted before each run. Separate code sharing nothing with the search
replays the complete $K=1$ domains ($4.3\times10^7$ and $3.5\times10^9$ states)
with identical survivors.

\paragraph{All 68 favourable CICYs, and a model the catalog misses} In 271
runs covering every (geometry, $|\Gamma|$) key of the catalog
($h^{1,1}\le6$, $K\le3$) the engine never misses a published model inside its
box (9,016 such models in the $K=2$ runs alone). It also keeps 80 orbits the catalog
does not list. The exact decision proves 53 of them unstable. Eighteen have
charge matrices of rank below $h-1$, where the decision does not apply, and are
unstable on numerical evidence only. The remaining nine are a single model,
seen on three CICY descriptions (5256, 5301, 5452) of one manifold. With
$|\Gamma|=4$ and charge vectors
$(-2,0,-1,1,0)$, $(0,-1,1,-2,1)$, $(0,2,1,0,-1)$, $(1,-2,1,0,0)$, $(1,1,-2,1,0)$,
it satisfies $c_1(V)=0$, $\ind V=-12$ and the anomaly and Bogomolov conditions,
and it is poly-stable with an exact witness: setting $t_5=1$, the coordinate
$t_4$ is the root of an irreducible degree-12 integer polynomial isolated in
$[642595/64736,\,769226/77493]$. The catalog has 2,128 models on this key, 656
of them in the same charge box, and none with this model's intersection
invariants. Its witness has $t_{\max}/t_{\min}\approx231$, so a stability check
confined to moderate Kähler ratios would not find it. It does not descend with
individually equivariant line bundles under any of the four free
$\Z_2\times\Z_2$ actions, so it is a new GUT model, not a new standard model.
Before trusting the verdict ``unstable'' we ran the decision on 548 published
models, all poly-stable. The first version rejected 91 of them; this exposed
two bugs, both fixed, and the check is now a regression test.

\paragraph{Downstairs spectra} Equivariant cohomology computed from the group
action reproduces all 27,821 published character tables. From our own survivor
set, 21 orbits on $\Xa/(\Z_2\times\Z_2)$ have three chiral families, no charged
chiral exotics and a Higgs doublet pair. All 10 published orbits are among them
and 11 are new. Each of the 11 lies outside the published scan's own stated
domain (its charge range, or its bound on vector-like pairs), which explains
why it is missing.

\paragraph{Superpotential} The flow formulation reproduces all thirteen operator
tables of the published standard-model database for its 2,012 complete models,
once nine unstated conventions are read off from the data. It also shows that
954 of the 2,007 published down-type Yukawa tables are internally inconsistent:
two families from the same line bundle receive different rows.

\paragraph{The physics verdict} We classified every configuration of the 21
orbits by explicit criteria: a nonzero top Yukawa at leading order, up- and
down-type Yukawa matrices that can reach rank three, and dimension-five proton
decay and dimension-four and bilinear $R$-parity violation forbidden to all
orders. No
orbit, published or new, meets them. Nine of the 11 new orbits fail on an exact
obstruction, as do seven of the 10 published ones; two new orbits (A04, A06)
remain undecided because the answer depends on higher-order coefficients and on
the singlet vacuum. None of the new orbits has a leading-order top Yukawa, nine
forbid dimension-five proton decay to all orders, and in one (A08) the
charge-allowed coupling $SLH_u$ vanishes cohomologically.

\paragraph{What is not claimed} No model is claimed to be the Standard Model.
Cohomology is that of a generic member of each family. Supersymmetry breaking,
moduli stabilisation and physical, canonically normalised Yukawa couplings are
not computed. That survivor sets stop changing beyond $K=4$ is recorded as an
observation, not as a finiteness theorem. Table~\ref{tab:results} lists the
main results with their status.

\begin{table}[h]
\centering\small
\setlength{\tabcolsep}{5pt}
\renewcommand{\arraystretch}{1.05}
\begin{tabularx}{\textwidth}{@{}X>{\raggedright}p{3.6cm}>{\raggedright\arraybackslash}p{4.1cm}@{}}
\toprule
result & status & domain \\
\midrule
GUT catalogs of $\Xa$ (10 and 58 models) and $\Xb$ (93) reproduced & exact set equality & each full catalog \\
complete $K=1$ domains replayed by independent code & exact & $4.3\times10^7$, $3.5\times10^9$ states \\
425 of 469 slices of $\Xa$ contain no model & theorem & all integer charges \\
$d_X\nmid3|\Gamma|$ excludes (7761, $|\Gamma|=2$) and (7447, $|\Gamma|=5$) & theorem & all integer charges \\
no published model inside a search box is missed & exact & 68 CICYs, 271 runs, $K\le3$ \\
poly-stable GUT model on CICY 5256 absent from the catalog & exact (algebraic witness) & $|\Gamma|=4$, $|k|\le2$ \\
53 of 80 extra survivors unstable; 18 more & exact; numerical only & charge-matrix rank $h-1$; below \\
27,821 published character tables reproduced & exact & whole database \\
21 downstairs orbits on $\Xa/(\Z_2\times\Z_2)$, 11 new, each absence explained & exact, generic member & upstairs $K=5$ survivors \\
13 operator tables reproduced; 954 of 2,007 $Y_d$ tables inconsistent & exact & 2,012 published models \\
no strong candidate; two orbits undecided & exact, perturbative & 21 orbits \\
survivor set of $\Xa$ unchanged for $K=4,\dots,7$ & observation, not a theorem & $\Gamma=\Z_2$ \\
\bottomrule
\end{tabularx}
\caption{Main results and the logical status of each.}
\label{tab:results}
\end{table}

\vspace{-6pt}
\paragraph{Open problems in the heterotic setting} Whether $\Xa$ carries
finitely many models is open: the anomaly bound on the charges degenerates as
the slope witness approaches the boundary of the Kähler cone, and we conjecture
that the index condition keeps it away. Couplings with Koszul level sum two or
more, which need a lift through the defining polynomial \cite{BBL17}, are
returned as unknown, and the 18 survivors of rank below $h-1$ need an exact
decision on positive-dimensional slope-zero loci. The survey should go to
larger $K$, where further models missing from the catalog would appear, and to
generalised CICYs \cite{AAGGL15} and hypersurfaces from the Kreuzer--Skarke
list \cite{KS00,AGHJN14}, where the Kähler cone is no longer the positive
orthant and the cone test must use the true one. An estimate of $10^{723}$
standard models on toric hypersurfaces \cite{CHL18} rules out box-by-box
exhaustion there, so the charge-unbounded statements become the useful ones.
Stability is decided at tree level; the one-loop correction \cite{BHW05} adds
the string coupling as one more variable for the same exact methods.

\newpage
\section{Outlook: the same methods in other compactifications}

What the engine exploits is a structure, not heterotic physics as such:
integer data; a supersymmetry condition that asks whether a cone meets a linear
space; a positive budget (anomaly or tadpole) that is spent along a branch of
the search; and abelian selection rules. The same structure appears in most
constructions that have produced the MSSM spectrum (Table~\ref{tab:outlook}).

\begin{table}[h]
\centering\footnotesize
\setlength{\tabcolsep}{4pt}
\renewcommand{\arraystretch}{1.1}
\begin{tabularx}{\textwidth}{@{}>{\raggedright}p{3.2cm}>{\raggedright}p{2.3cm}>{\raggedright}p{3.0cm}>{\raggedright}X>{\raggedright\arraybackslash}p{2.9cm}@{}}
\toprule
construction & integer data & budget & supersymmetry condition & selection rules \\
\midrule
heterotic line bundles (here) & charge matrix & $c_2(TX)-c_2(V)\ge0$ & slopes, quadratic in $t$; one loop \cite{BHW05} & $U(1)^4$ (flows) \\
IIB, magnetised D7-branes & divisors, fluxes & D7 and D3 tadpoles & FI terms, linear in $t$ & anomalous $U(1)$s, E3 \\
IIB closed-string fluxes & $F_3$, $H_3$ & $N_{\rm flux}\le Q_{\rm D3}$ & $DW=0$, polynomial near LCS & --- \\
F-theory & $G_4$ flux & $\chi(Y)/24-\tfrac12G_4^2\ge0$ & primitivity, linear in $J$ & $U(1)$s, discrete \\
IIA, D6 on $T^6/\Z_2\times\Z_2'$ & wrapping numbers & RR tadpoles, K-theory & linear in complex-structure functions & anomalous $U(1)$s \\
\bottomrule
\end{tabularx}
\caption{Where the ingredients of the search reappear (LCS: large complex structure).}
\label{tab:outlook}
\end{table}

\vspace{-4pt}
\paragraph{Type IIB with magnetised D7-branes} A $U(1)$ stack on a divisor
$D_a$ with worldvolume flux $F_a$ has, at leading order, Fayet--Iliopoulos term
$\int_XJ\wedge D_a\wedge F_a$, \emph{linear} in the Kähler class. Whether all
of them vanish at one point of the Kähler cone is then the question of
Section~2, whether a rational subspace meets an open cone, but now necessary
\emph{and} sufficient, and settled by linear programming with an exact Farkas
certificate either way. Chiral indices are integer intersection numbers, the D7
tadpole is a linear equality, Freed--Witten quantisation is a parity condition,
and the D3 tadpole is a budget that worldvolume and closed-string fluxes share.
The $SU(5)$ models of \cite{BBGW08} show how restrictive these conditions are.
A first experiment: exhaust the three-family $SU(5)$ D7 configurations on
orientifolds of small-$h^{1,1}$ CICYs, and apply the flow rule to the anomalous
$U(1)$s with E3-instanton charges added as generators, to decide which
perturbatively forbidden couplings, the $\mathbf{10}\,\mathbf{10}\,\mathbf5$
top coupling among them, an instanton could supply.

\paragraph{Flux compactification and the shared tadpole} Fluxes fix the complex
structure and the dilaton \cite{GKP02}, and the Kähler moduli need
non-perturbative effects \cite{KKLT03} or $\alpha'$ corrections
\cite{BBCQ05}. The tadpole conjecture \cite{BBGL20} states that
fixing all complex-structure moduli costs more D3 charge than is available when
there are many moduli. Two computations follow directly. The perturbatively flat
vacua of \cite{DKMM20} are pairs of integer flux vectors such that a derived
vector lies in the Kähler cone of the mirror and is orthogonal to one of the
flux vectors, subject to integrality conditions and the tadpole bound: a cone
meeting a subspace, and a monotone budget, which is what our search is built
for. A complete, certified list of
such vacua on small-$h^{2,1}$ examples would replace sampling with a theorem.
The second asks how much of the D3 tadpole the cheapest certified three-family
brane sector uses on a given orientifold, and whether the rest can fix the
complex structure; this tests the conjecture example by example as a finite
computation. Visible sectors on D3-branes at del Pezzo singularities
\cite{CKMQV12}, which meet the global data mainly through the tadpole, are a
natural first target.

\paragraph{F-theory} In the $\mathcal O(10^{15})$ compactifications with the
exact chiral spectrum of the Standard Model \cite{CHLLT19}, global consistency
reduces to a criterion on the base; the vector-like spectrum and the Yukawa
couplings must still be settled model by model. Proving vector-like exotics
absent needs upper bounds on cohomology, that is lower bounds on ranks, which a
minor modulo a prime certifies, and the flow rule applies unchanged to
Mordell--Weil $U(1)$s. In $G_4$ searches the chiral indices are linear in the
flux, and since the intersection form is positive definite on primitive
$(2,2)$-classes the D3 tadpole confines the flux to an ellipsoid, which can be
enumerated exactly, branch by branch.

\paragraph{Type IIA intersecting D6-branes} On $T^6/(\Z_2\times\Z_2')$ the
supersymmetry equality of each stack is linear in four positive functions of
the complex-structure moduli \cite{CSU01,BCLS05}. A common supersymmetric point
therefore needs a positive vector orthogonal to the span of the stacks' wrapping
data: the slice problem of Section~2, decided by the same Gordan test and
monotone as stacks are added. All vacua of this orientifold have been counted
exactly \cite{LS22}, and MSSM-like ones estimated at one in a billion
\cite{GBHLW06}; restricting the exact count to the MSSM chiral spectrum and
running the flow analysis on the anomalous $U(1)$s would give a certified list
of three-family models together with the couplings each allows perturbatively.

\smallskip
In each case the work has the same shape: find the integer data, state every
condition as exact, necessary or a declared restriction, and apply the monotone
ones as early in the tree as possible. A search built this way settles what
sampling cannot: that a model is absent, that a catalog is complete in a box,
or that an allowed coupling vanishes.

\vspace{-2pt}
\noindent\rule{\textwidth}{0.4pt}
\vspace{-14pt}
{\scriptsize
\renewcommand{\section}[2]{}
\begin{multicols}{2}
\begin{thebibliography}{99}
\setlength{\itemsep}{0pt}\setlength{\parsep}{0pt}
\bibitem{AGLP11} L.~B.~Anderson et al., Phys.\ Rev.\ D\,84 (2011) 106005 [1106.4804].
\bibitem{AGLP12} L.~B.~Anderson et al., JHEP\,06 (2012) 113 [1202.1757].
\bibitem{ACGLP13} L.~B.~Anderson et al., JHEP\,01 (2014) 047 [1307.4787].
\bibitem{BBCL16} S.~Blesneag et al., JHEP\,01 (2016) 152 [1512.05322].
\bibitem{BBL17} S.~Blesneag et al., JHEP\,01 (2017) 119 [1607.03461].
\bibitem{AAGGL15} L.~B.~Anderson et al., arXiv:1507.03235.
\bibitem{KS00} M.~Kreuzer, H.~Skarke, Adv.\ Theor.\ Math.\ Phys.\,4 (2002) 1209 [hep-th/0002240].
\bibitem{AGHJN14} R.~Altman et al., arXiv:1411.1418.
\bibitem{CHL18} A.~Constantin et al., arXiv:1810.00444.
\bibitem{BHW05} R.~Blumenhagen et al., JHEP\,06 (2005) 020 [hep-th/0504232].
\bibitem{BBGW08} R.~Blumenhagen et al., Nucl.\ Phys.\ B\,815 (2009) 1 [0811.2936].
\bibitem{GKP02} S.~B.~Giddings et al., Phys.\ Rev.\ D\,66 (2002) 106006 [hep-th/0105097].
\bibitem{KKLT03} S.~Kachru et al., Phys.\ Rev.\ D\,68 (2003) 046005 [hep-th/0301240].
\bibitem{BBCQ05} V.~Balasubramanian et al., JHEP\,03 (2005) 007 [hep-th/0502058].
\bibitem{BBGL20} I.~Bena et al., arXiv:2010.10519.
\bibitem{DKMM20} M.~Demirtas et al., arXiv:1912.10047.
\bibitem{CKMQV12} M.~Cicoli et al., arXiv:1206.5237.
\bibitem{CHLLT19} M.~Cveti\v{c} et al., Phys.\ Rev.\ Lett.\,123 (2019) 101601 [1903.00009].
\bibitem{CSU01} M.~Cveti\v{c} et al., Phys.\ Rev.\ Lett.\,87 (2001) 201801 [hep-th/0107143].
\bibitem{BCLS05} R.~Blumenhagen et al., Ann.\ Rev.\ Nucl.\ Part.\ Sci.\,55 (2005) 71 [hep-th/0502005].
\bibitem{LS22} G.~J.~Loges, G.~Shiu, arXiv:2206.03506.
\bibitem{GBHLW06} F.~Gmeiner et al., JHEP\,01 (2006) 004 [hep-th/0510170].
\end{thebibliography}
\end{multicols}}

\end{document}
